% Living discrepancy ledger for the four primary DKPS theorem targets.
% Update this file whenever a public theorem changes its assumptions,
% conclusion, estimator, asymptotic regime, or relationship to the source paper.

\documentclass[11pt]{article}

\usepackage[T1]{fontenc}
\usepackage{lmodern}
\usepackage{microtype}
\usepackage[margin=1.0in]{geometry}
\usepackage{amsmath,amssymb,mathtools}
\usepackage{booktabs,longtable,array}
\usepackage[colorlinks=true,linkcolor=black,urlcolor=blue,citecolor=blue]{hyperref}

\newcommand{\Lean}{\textsc{Lean}}
\newcommand{\DKPS}{\textsc{DKPS}}
\newcommand{\whp}{\text{with high probability}}
\newcommand{\norm}[1]{\left\lVert #1\right\rVert}
\newcommand{\code}[1]{\texttt{#1}}
\setlength{\tabcolsep}{4pt}
\setlength{\emergencystretch}{2em}

\title{The Four Main DKPS Papers:\\Differences Between the Literature and the Lean Formalization}
\author{Jonathan Crall \and GPT-5.6 Thinking}
\date{\today}

\begin{document}
\maketitle

\begin{abstract}
This is a living audit of the four primary paper targets in the
\DKPS{} formalization repository.  Its purpose is not to restate every theorem,
but to record every material difference between the published mathematical
presentation and the theorem that is actually proved in \Lean.  Differences are
classified as repaired statements, surfaced assumptions, strengthened or
weakened conclusions, modelling choices, and still-open composition gaps.  The
source of truth remains the theorem signatures and proofs; this document is the
human-readable discrepancy ledger that should change with them.
\end{abstract}

\section{Scope and maintenance rule}

The four direct targets are:
\begin{enumerate}
\item Acharyya, Trosset, Priebe, and Helm (2024), consistent estimation of
      generative-model representations in the data kernel perspective space;
\item Acharyya, Agterberg, Park, and Priebe (2025), concentration bounds for
      response-based vector embeddings;
\item Helm, Acharyya, Park, Duderstadt, and Priebe (2025), statistical inference
      on black-box generative models in DKPS;
\item Helm, Johnson, and Priebe (2026), query-efficient model evaluation using
      cached responses (Quench).
\end{enumerate}

A row belongs in this document when any of the following changes:
\begin{itemize}
\item an assumption present in \Lean{} is absent or only implicit in the paper;
\item a paper statement is false or underspecified without a repair;
\item the formalization proves a stronger or weaker conclusion;
\item the formalization uses a different estimator or MDS variant;
\item a probability, measurability, compactness, or selection issue is made
      explicit;
\item a formal theorem covers only a special case of the paper, or extends it.
\end{itemize}

Compatibility theorems may retain older, heavier interfaces.  The preferred
public capstone is the theorem with the smallest honest hypothesis load.  A
proof refactor that merely shortens an argument does not change this ledger; a
refactor that removes a caller-visible condition does.

\section{Acharyya et al. (2024): asymptotic DKPS and raw-stress MDS}

\begin{longtable}{>{\raggedright\arraybackslash}p{0.18\textwidth}
                  >{\raggedright\arraybackslash}p{0.25\textwidth}
                  >{\raggedright\arraybackslash}p{0.36\textwidth}
                  >{\raggedright\arraybackslash}p{0.13\textwidth}}
\toprule
Topic & Literature & Formalization & Classification \\
\midrule
\endhead
Choice of population minimizer & The fixed-configuration theorem is phrased
up to affine ambiguity, with a subsequence selecting a compatible limit. & The
unconditional theorem is set-valued: estimated minimizers approach the set of
population pair-distance profiles.  A literal fixed-configuration theorem
requires the explicit predicate \code{UniquePairProfile}. & Repaired statement
\\
\addlinespace
Subsequence versus full sequence & The paper uses subsequences and a diagonal
argument. & Once uniqueness of the pair-distance profile is assumed, the
formalization obtains full-sequence convergence; the diagonal argument is no
longer needed. & Stronger conclusion under explicit uniqueness \\
\addlinespace
Measurable selection & A measurable choice of MDS minimizer is left implicit. &
The proof uses a deterministic modulus of continuity and event inclusion, so no
measurable selector is required. & Assumption removed \\
\addlinespace
Configuration error & Informal statements identify configurations only up to
affine transformation. & \code{ConfigError} is direct coordinate error and is
therefore stronger than affine-quotiented closeness.  The repaired consistency
layer primarily compares pairwise-distance profiles where that is the invariant
quantity. & Modelling distinction \\
\addlinespace
Moment assumptions & The sampling discussion is iid. & The variance algebra is
proved under pairwise independence, which is sufficient.  Measurability,
integrability, and \(L^2\) hypotheses are explicit. & Generalization plus
surfaced technical assumptions \\
\addlinespace
Distribution-over-model conclusion & The paper also discusses an
\(L^p\)-over-model-distribution form. & That form is not yet formalized in the
2024 library. & Scope gap \\
\bottomrule
\end{longtable}

The important repair is conceptual: nonuniqueness of raw-stress MDS is treated
as a property of the theorem, not hidden inside a choice function.  The
formalization's strongest unconditional result is therefore about the minimizer
set or its invariant pair-distance profile.  A fixed representative is available
only after a uniqueness condition makes it mathematically meaningful.


\paragraph{Role in the Perfect Quench scaffold.}
The new raw-response track reuses the 2024 matrix-valued sample-mean
second-moment theorem rather than postulating response means.  It adds an
independent product probability space for reference and response randomness,
then lifts the existing replicate-average calculation through random reference
selection.  This is composition work beyond the paper, not a change to the
2024 consistency theorem itself.

\section{Acharyya et al. (2025): finite-sample response embeddings}

\begin{longtable}{>{\raggedright\arraybackslash}p{0.18\textwidth}
                  >{\raggedright\arraybackslash}p{0.25\textwidth}
                  >{\raggedright\arraybackslash}p{0.36\textwidth}
                  >{\raggedright\arraybackslash}p{0.13\textwidth}}
\toprule
Topic & Literature & Formalization & Classification \\
\midrule
\endhead
Population spectral assumptions & Rank \(d\), positive spectrum, and upper
spectral control are stated at paper level. & Positive semidefiniteness, rank,
and eigenvalue floor/cap are represented as exact matrix predicates.  In the
preferred Quench-facing capstones, a population Gram witness now derives
positive semidefiniteness and rank automatically. & Surfaced and partially
discharged \\
\addlinespace
Smallness regime & The paper packages the regime through asymptotic growth such
as sufficiently many responses relative to the model count. & The finite
spectral theorem exposes the exact entrywise, eigengap, and polar-factor
smallness inequalities consumed by the proof.  \code{GrowingConfigControl}
packages their eventual validity for varying matrix size. & Quantifiers made
explicit \\
\addlinespace
Alignment choice & The paper chooses an optimal orthogonal alignment
classically. & The formalization separates the existential alignment predicate
from a noncomputable selected alignment, so probability arguments do not depend
on measurability of a data-dependent choice. & Selection issue resolved \\
\addlinespace
Rates & The paper presents sharper asymptotic bookkeeping. &
\code{cmdsEntrywiseRate}, \code{configBound}, and the composed end-to-end rate
are valid but intentionally loose.  No claim of optimal constants or exponents
is made. & Weaker quantitative sharpness \\
\addlinespace
Growing target-augmented CMDS & The paper's theorem is presented for a fixed
finite model collection. & The formalization additionally supports a growing
\(n+1\) batch containing the \(n\) sampled references and the current target,
and proves pairwise-distance control without a globally aligned coordinate
map. & Extension beyond paper \\
\addlinespace
Response boundedness & Bounded dissimilarities are used in the concentration
pipeline. & The preferred response bridge needs only a population response-norm
envelope.  On the response-mean event, the random sample norm and both sample
and population dissimilarity bounds are derived internally. & Hypotheses
reduced \\
\addlinespace
Davis--Kahan implementation & The paper invokes spectral perturbation theory. &
The production proof currently uses a proved bespoke cross-energy and polar
factor chain.  Newer general Davis--Kahan results are intentionally deferred
unless they remove a caller-visible condition, such as the polar smallness
hypothesis. & Deliberate implementation choice \\
\bottomrule
\end{longtable}

The main remaining paper-facing burden is the uniform lower spectral bound.  It
is a genuine identifiability condition for the spectral estimator rather than a
proof-engineering artifact.  The upper spectral envelope and the eventual rate
certificate are also still supplied uniformly over stages and targets.


\paragraph{Role in the Perfect Quench scaffold.}
The completion program replaces stagewise global spectral assumptions by a
high-probability certificate derived from a population covariance floor and iid
reference sampling.  It also supplies an explicit conservative response
schedule.  These changes reduce the assumptions of the Quench composition; the
proved 2025 concentration statements remain intact until the new event-level
spectral bridge is completed.

\section{Helm et al. (2025): statistical inference transfer}

\begin{longtable}{>{\raggedright\arraybackslash}p{0.18\textwidth}
                  >{\raggedright\arraybackslash}p{0.25\textwidth}
                  >{\raggedright\arraybackslash}p{0.36\textwidth}
                  >{\raggedright\arraybackslash}p{0.13\textwidth}}
\toprule
Topic & Literature & Formalization & Classification \\
\midrule
\endhead
Continuity assumptions & Assumptions A2 and A4 are pointwise in the perturbed
coordinate. & The proved transfer theorem packages observations as
embedding--label pairs and uses joint continuity of the learning rule and loss.
Both the paper-shaped predicates and the stronger predicates used by the proof
are exposed. & Stronger sufficient assumptions \\
\addlinespace
Compactness assumption A3 & The paper states closedness, boundedness, and
completeness of the image. & The main proof uses a uniform compact-range
condition, equivalent in finite-dimensional Euclidean space; a literal
paper-shaped predicate is also retained. & Equivalent finite-dimensional
encoding \\
\addlinespace
Measurability & Mostly implicit. & Measurability of estimated embeddings is
explicit because risks and expectations are taken. & Surfaced technical
assumption \\
\addlinespace
Maximum and budgets & Equation (3) uses a finite maximum and two asymptotic
budgets. & The formalization uses a finite supremum and one abstract budget
index, which may encode a joint schedule. & Representation change \\
\addlinespace
MDS estimator used by the bridge & The argument cites the eigengap-free
asymptotic raw-stress consistency result from the 2024 paper. & The current
end-to-end bridge realizes the practical estimator as classical spectral MDS
and therefore requires a uniform lower population eigenvalue bound.  This
assumption is not stated in the Helm paper. & Material surfaced discrepancy \\
\bottomrule
\end{longtable}

The last row is the most important discrepancy in this paper.  There are two
honest ways to close it: connect Helm directly to the raw-stress consistency
route used in the paper, or amend the spectral-estimator theorem with the
necessary eigenvalue-stability condition.  The current formalization takes the
second route and states the condition explicitly.


\paragraph{Role in the Perfect Quench scaffold.}
No new inference-transfer theorem is required for the fixed-subset Quench
result.  The reusable contribution from this paper track is its discipline
around continuity, compactness, and measurable event composition.  The ledger
therefore records no artificial Helm-specific dependency; future reuse should
be added only when it removes a Quench-facing hypothesis.

\section{Helm, Johnson, and Priebe (2026): Quench}

\begin{longtable}{>{\raggedright\arraybackslash}p{0.18\textwidth}
                  >{\raggedright\arraybackslash}p{0.25\textwidth}
                  >{\raggedright\arraybackslash}p{0.36\textwidth}
                  >{\raggedright\arraybackslash}p{0.13\textwidth}}
\toprule
Topic & Literature & Formalization & Classification \\
\midrule
\endhead
Nearest-neighbor ties & The estimator is described as averaging cached scores
of nearest references. & The preferred theorem uses the literal tie-averaged
estimator.  An older selected-neighbor compatibility theorem is retained but is
not the strongest paper-facing result. & Faithfulness repair \\
\addlinespace
Coverage assumption & The paper describes positive probability on relevant
compact regions. & \code{PerspectiveFullSupport} requires positive mass in
every positive-radius ball centered on the perspective range, exactly what the
finite-net iid coverage proof consumes. & Measure-theoretic repair \\
\addlinespace
Reference sampling & Iid sampling is described informally. &
\code{IIDReferenceSampler} records the finite joint-product law needed for the
miss-probability calculation.  Coverage is then proved rather than assumed. &
Surfaced probabilistic interface \\
\addlinespace
Finite factorization of the model class & Early formal capstones represented all
models in one fixed finite configuration. & The strongest theorem adjoins the
current target to the growing reference sample and runs CMDS on
\(\operatorname{Fin}(n+1)\).  No global model-to-finite-index map or globally
aligned estimated perspective is required. & Assumption removed \\
\addlinespace
Population CMDS validity & The paper treats the population embedding as given. &
A supplied population Gram witness now derives positive semidefiniteness and
rank at most \(d\); callers no longer prove those consequences separately. &
Hypotheses reduced \\
\addlinespace
Dissimilarity boundedness & Bounded sample and population dissimilarities were
previously separate bridge assumptions. & The preferred finite-model capstone
requires one population response-norm envelope.  Sample boundedness is derived
only on the high-probability response-mean event, where it is actually needed. &
Hypotheses reduced \\
\addlinespace
Compactness for finite model classes & Previously explicit. & Finiteness of the
model class derives compactness of the perspective range. & Assumption removed
\\
\addlinespace
Uniform response concentration & The paper's broad model-class presentation
suggests uniform control. & The complete second-moment derivation currently
uses a finite target class and a double union bound.  For an infinite class,
uniform concentration remains an explicit input. & Major remaining scope gap
\\
\addlinespace
All query subsets & The paper's efficiency claim ranges over query budgets. &
The coordinate-level library has fixed-subset, size-\(m\), and all-proper-subset
capstones through per-subset certificates.  The strongest growing response
bridge currently closes one fixed subset at a time; a uniform family of
spectral and Lipschitz certificates is still application-level input. &
Composition gap \\
\bottomrule
\end{longtable}

\section{Perfect Quench completion scaffold}

The repository now contains an executable proof scaffold under
\code{DkpsQuench/Perfect}.  Its 61 open proof bodies are split so the count
reflects the amount of mathematical work rather than hiding the program in one
final declaration.

\begin{longtable}{>{\raggedright\arraybackslash}p{0.22\textwidth}
                  >{\raggedright\arraybackslash}p{0.10\textwidth}
                  >{\raggedright\arraybackslash}p{0.30\textwidth}
                  >{\raggedright\arraybackslash}p{0.29\textwidth}}
\toprule
Track & Open proofs & New mathematical content & Hypotheses removed from the current bridge \\
\midrule
\endhead
Population geometry & 4 & Centering, exact Euclidean double centering, and the
augmented radial identity. & Separate configuration, Gram, radial, PSD, and rank
inputs collapse to one response-distance realization. \\
\addlinespace
Spectral regularity & 15 & Coordinate mean/product weak laws, covariance algebra,
scalar-event measurability, fixed-dimensional
covariance weak law, entrywise-to-quadratic transfer, reference scatter, target
augmentation, and a trace ceiling. & Global floor and ceiling assumptions over
every reference outcome are replaced by one population covariance
nondegeneracy condition and a measurable high-probability event. \\
\addlinespace
Raw responses & 9 & Independent product probability space, concrete replicate
averages, second moments after random model selection, and a finite-class union
bound. & Abstract response means, manual moment events, and per-index
measurability obligations. \\
\addlinespace
Compactness and regularity & 8 & Finite and compact response envelopes,
pathwise-raw to sample/population Lipschitz bridges, and polynomial Euclidean
covering numbers. & Explicit population bounds, selected finite nets, entropy
exponents, covering constants, and separately supplied regularity certificates.
\\
\addlinespace
Infinite-class concentration & 8 & Finite stage nets, measurable finite
subevents, entropy union bounds, and deterministic Lipschitz extension. & The
finite-model restriction or an assumed uniform concentration event. \\
\addlinespace
Explicit schedules & 10 & Conservative polynomial tolerance, canonical
shrinking-net radius, entropy-aware replicate budgets, and all limit
calculations exposed. & Caller-selected nets, entry rates, smallness
inequalities, polar bookkeeping, and the full
\code{GrowingConfigControl} certificate. \\
\addlinespace
Spectral composition & 3 & Intersection of response and spectral subevents with
the existing proved pairwise CMDS perturbation theorem. & Global spectral facts
inside the response bridge. \\
\addlinespace
Final composition & 4 & Finite and compact-infinite fixed-subset capstones plus
the two all-proper-subset lifts. & All intermediate certificates from the public
Perfect Quench interface. \\
\bottomrule
\end{longtable}

The four final theorem statements already expose the intended endpoint.  The
finite theorem assumes raw iid response replicates, exact population
response-distance realization, a population covariance floor, full support,
and score Lipschitzness.  The compact infinite-model theorem adds only compact
perspective range and one pathwise raw-response Lipschitz condition; polynomial
nets, entropy growth, replicate-mean/population-mean regularity, and response
norm envelopes are derived internally.  The all-query theorems
allow each proper query subset to carry its own perspective and response model
while sharing the reference and response probability spaces.

The production cross-energy and polar-factor proof is deliberately retained.
General Davis--Kahan results should replace it only if they eliminate or weaken
a caller-visible assumption, especially the polar smallness condition.  A
constant improvement or shorter proof alone is not part of this completion
program.

\section{Current hypothesis-reduction frontier}

The scaffold defines the remaining frontier precisely:
\begin{enumerate}
\item complete the exact population double-centering identities;
\item derive high-probability spectral floor and ceiling events from compact iid
      references and population covariance nondegeneracy;
\item connect raw iid response arrays to the augmented response event on an
      independent product probability space;
\item derive population response envelopes, raw-to-mean regularity, and
      polynomial finite nets from finiteness or compactness;
\item derive infinite-class uniform concentration from those finite nets;
\item discharge the conservative polynomial schedule and event-level spectral
      composition;
\item assemble the finite, compact-infinite, and all-query capstones.
\end{enumerate}

No cleaner-only spectral refactor is on this list.  Once these declarations are
proved, remaining work should be rate sharpening or assumption comparison, not
construction of another theorem layer.

\section{Audit checklist for future edits}

For every new preferred capstone, record:
\begin{enumerate}
\item the exact paper theorem or algorithm step it represents;
\item each assumption not written in the source paper;
\item each paper assumption discharged internally by a preceding theorem;
\item whether the estimator is raw-stress MDS, classical MDS, or another
      explicitly named variant;
\item whether the conclusion is fixed-target, uniform over targets, fixed-query,
      or uniform over query subsets;
\item whether rates are qualitative, explicit but loose, or quantitatively
      matched to the paper;
\item any finite-dimensional, finite-model, measurability, separability, or
      compactness restriction.
\end{enumerate}

\begin{thebibliography}{9}
\bibitem{acharyya2024}
A. Acharyya, M. W. Trosset, C. E. Priebe, and H. S. Helm.
\newblock Consistent estimation of generative model representations in the data
kernel perspective space.
\newblock arXiv:2409.17308, 2024.

\bibitem{acharyya2025}
A. Acharyya, J. Agterberg, Y. Park, and C. E. Priebe.
\newblock Concentration bounds on response-based vector embeddings of black-box
generative models.
\newblock arXiv:2511.08307, 2025.

\bibitem{helm2025}
H. S. Helm, A. Acharyya, Y. Park, B. Duderstadt, and C. E. Priebe.
\newblock Statistical inference on black-box generative models in the data
kernel perspective space.
\newblock Findings of ACL, 2025.

\bibitem{quench2026}
H. Helm, B. Johnson, and C. E. Priebe.
\newblock Query-efficient model evaluation using cached responses.
\newblock arXiv:2605.07096, 2026.
\end{thebibliography}

\end{document}
